%\vspace{-3mm}
\section{Introduction} \label{sec:introduction}
%\vspace{-3mm}

% \sg{
% Start directly saying that in Legal Analytics especially in Common Law countries, there are two practical and popular tasks - given a legal case, identify relevant statutes and precedents (prior cases). \\
% There have been many works on both tasks, but every prior work has focused on only one task. In reality, Law practitioners often have to do both tasks together. But till date, there is even no dataset that enables development of models for doing both tasks together. The objective of this work is to bridge this gap. \\
% We develop the first dataset for ... Describe the dataset we develop briefly (point to the dataset comparison table).\\
% We benchmark a wide variety of methods on the dataset ... state some methods.\\
% This first-of-its-kind dataset and analyses bring out some surprising observations - lexical methods perform well in one task while semantic methods perform well in the other. A probable reason for this surprising observation is the difference in the nature of the two tasks (explain). For this reason, multi-tasking approaches also do not achieve much improvement.\\
% Finally we develop simple ensemble approaches that combine the two types of methods, and successfully achieve improvements over both tasks.\\
% We also perform experiments over datasets from other legal jurisdictions and verify that the ensemble approaches achieve improvement. \\
% End with a recap of our main contributions.
% }\paul{Addressed}
%These can be time consuming activities and may suffer from recall issues. 

\noindent In the legal domain, laws (statutes) and prior cases are considered to be the fundamental sources of knowledge that guide principles of jurisdiction~\cite{joshi2023ucreat}. 
In practice, a legal practitioner when faced with a legal case, typically uses their experience and knowledge to identify prior precedents and applicable statutes in the given situation. 
This process is time-consuming, and the problem worsens with the growing number of legal cases in populous countries like India \cite{malik2021ildc}. Hence, there is an imminent need to automate the process to make it more efficient. 

% \noindent Two tasks have been proposed independently in this regard: Legal Statute Retrieval (LSR) -- that aims to identify/retrieve statutes that are applicable in a given query case \cite{paul2024legal}, and Prior Case Retrieval (PCR) -- that aims to identify/retrieve relevant prior cases/precedents that should be cited by the given query case \cite{joshi2023ucreat}. Traditionally, the two tasks have been modeled separately leading to creation of different corpora and models for each of them. However, both the tasks are inherently related. For example, a legal case would typically cite another prior legal case if they share similar factual situation, consequently both would be citing similar statutes. For computational models to exploit such relationships a corpus which combines both the tasks is required. We address this gap by making following contributions:

\noindent Two tasks have been proposed independently in this regard: Legal Statute Retrieval (LSR) -- that aims to identify/retrieve statutes that are applicable in a given query case \cite{paul2024legal}, and Prior Case Retrieval (PCR) -- that aims to identify/retrieve relevant prior cases/precedents that should be cited by the given query case \cite{joshi2023ucreat}. Traditionally, the two tasks have been modeled separately leading to creation of different corpora and models for each of them. However, both tasks are inherently connected in real legal practice. In India, as well as in many other countries, case law is treated as a primary source of law supplementing statutes, regulations, and constitutional provisions~\citep{manupatra-guide}. 
%and official legal infrastructures such as the India Code Information System (ICIS: \url{https://www.indiacode.nic.in}) provides mechanisms to locate the relevant precedents and amendments while retrieving statutes.
For example, a legal case would typically cite another prior legal case if they share similar factual situation, and consequently both would be citing similar statutes. For computational models to exploit such connections, a corpus which combines both the tasks is required. We address this gap by making the following contributions:



%It makes it difficult to compare models across tasks, since the models were not aware in any way about the other task during training. \paul{However, both the tasks are essentially legal document retrieval tasks, hence it would be interesting to explore if it is possible to utilize the outputs/training gained by one task in a bid to improve the other.} This requires a common corpus for both the tasks having query cases (with missing precedent citations and missing statutes) and corresponding pool of candidate prior cases and pool of candidate statutes. We address this gap in this paper. In a nutshell, we make the following contributions:
%statues and prior precedents are correlated. For example, a case citing a prior precedent is likely to share similar factual situation and consequently both the cases are likely to invoke similar statues. In this work, explore the possibility of modeling both the tasks simultaneously. However, for solving both the tasks together there is a lack of any corpus having query cases (with missing precedent citations and missing statutes) and corresponding pool of candidate prior cases and pool of candidate statutes. In this scenario, we make the following contributions:
\begin{itemize}[noitemsep,nosep,leftmargin=*]
\item  In this paper, we propose \ilpcsr\ (Indian Legal Corpus for Prior Case and Statute Retrieval), a large corpus of query cases along with a candidate pool of prior cases and statutes in English for the Indian legal domain. To the best of our knowledge, we are first to develop such a corpus having both prior cases and statutes for the same set of queries. 
\item We perform extensive experiments with a variety of models for  each task, including lexical models (e.g., BM-25) and semantic models (e.g., transformer-based models and GNNs). Further, we experiment with an ensemble of semantic and lexical models, which improves performance in both tasks.
%\paul{Finally, we develop a prompting-based technique with an LLM to re-rank the top candidates given by the ensemble model.}
\item To test our intuition about inherent dependencies between the two tasks, we develop two methods: 
(i)~joint prediction model for both tasks via multi-task learning model based on GNNs; 
(ii)~a unique pipeline based approach comprising of an ensemble model that acts as a retriever, and an LLM as a re-ranker that performs tasks in a sequential order. The pipeline-based approach (in either order LSR $\rightarrow$ PCR or PCR $\rightarrow$ LSR) gives better performance than lexical/semantic models on each task separately as well as the joint multi-task model.

\item In-depth analyses on the corpus and experiments brings out some interesting observations, for example, lexical methods perform well in one task (PCR) while semantic methods perform well in the other (LSR). 
This may be due to differences in candidate types. Statutes are short, technical, and abstract, while precedents are longer and linguistically closer to the query, potentially explaining the better performance of the lexical-semantic ensemble model.
We release \ilpcsr\ and developed models under open-source license via Hugging Face~\footnote{\url{https://huggingface.co/datasets/Exploration-Lab/IL-PCSR}} and GitHub~\footnote{\url{https://github.com/Exploration-Lab/IL-PCSR}}. %upon acceptance. For now, we make a small sample of the dataset available anonymously.\footnote{ \url{https://anonymous.4open.science/r/sample_dataset-F620/README.md}}
%Thereby, an Ensemble-based model that effectively combines lexical and semantic features is used for both tasks in the retrieval stage. Further, we note that fine-tuning both tasks in a pipeline fashion led to the best performance, compared to independent or multi-task fine-tuning. 
%\item \paul{Experiments and analyses on the corpus brings out some interesting observations -- lexical methods perform well in one task (PCR) while semantic methods perform well in the other (LSR). A probable reason is due to differences in the nature of statute candidates, which are usually short, technical and abstract, as compared to precedent candidates which are long, and have a similar language as the query. Thereby, an Ensemble-based model that effectively combines lexical and semantic features is used for both tasks in the retrieval stage. Further, we note that fine-tuning both tasks in a pipeline fashion led to the best performance, compared to independent or multi-task fine-tuning.}
%\item \paul{Following on from such an insight, we constructed an effective prompt template for GPT-4.1, where we first perform LSR independently, and then provide the statutes predicted as relevant as part of the prompt for PCR. Conversely, the ground-truth statutes of predicted precedents are added to the re-ranking pool for LSR to improve the performance. This establishes our hyothesis further --- utilizing one task effectively can lead to improvements for the other task.}
% \item Experiments and analyses on the corpus brings out some interesting observations -- lexical methods perform well in one task (PCR) while semantic methods perform well in the other (LSR). A probable reason for this surprising observation is the difference in the nature of the two tasks, such as, statutes are short and abstract, and the overall semantics of the query help to identify the relevant statutes. 
% However, precedents are long and have similar language as the query (allows better lexical matches). 
% Moreover, different portions of the query text (that have different local semantics) match accurately with different precedents. For these reasons, multi-task approaches (solving both the tasks simultaneously) also do not show much improvement (\S\ref{sec:expt}).
% \item Based on experimental observations, we propose an ensemble-based approach  for combining scores coming from semantic and lexical methods respectively. These techniques use a linear interpolation of semantic and lexical scores, with a parameter $\alpha$ controlling the weightage between the two scores. We derive the $\alpha$ value by grid-search. Using this simple ensemble technique, we are able to obtain gains over both the individual methods. This observation holds true for both statute and precedent retrieval, across different combinations of lexical and semantic approaches.  
% \item We further experiment with datasets from other jurisdictions. We experiment with the \coliee{} datasets, specifically \texttt{COLIEE 2024 Task 3} (Statute Retrieval on Japanese law) and \texttt{COLIEE 2024 Task 1} (Precedent Retrieval on Canadian Federal cases). We verify that the same trends as \ilpcsr{} hold for \coliee{}, and that the ensemble approach provides an improvement here as well.  
%\item The \ilpcsr\ corpus, along with all model implementations, will be released publicly upon acceptance. For now, we make a small sample of the dataset available anonymously.\footnote{ \url{https://anonymous.4open.science/r/AnonymousSampleDataset-67B6/readme.md}}
%We will release the dataset and model code upon acceptance. 
\end{itemize}

 

%Most countries across the world follow either the civil law system (cases are decided on the basis of statutes, or well-defined laws) or common law system (cases are decided on the basis of decisions and observations of similar prior cases). In countries like India, which follow a mixture of civil and common law, identifying the relevant statutes and precedents given the description of legal facts or a legal situation thus becomes one of the most studied tasks in legal NLP. Also, since India has a huge disparity in the ratio of number of legal professionals to the general population, as well as a backlog of several millions of cases, automation of these two tasks can lead to improved efficiency and speed in the judicial system.

%From early efforts involving mathematical and statistical methods, to using lexical and hand-crafted features, to pre-trained embeddings and transformer-based models lately, there has been continuous effort to improve both tasks of legal statute retrieval (LSR) and prior case retrieval (PCR). While LSR has usually been modeled as a multi-label/single-label classification task (considering a limited number of statutes), PCR has always been modeled as the classic Information Retrieval problem. While there are many prior works on both LSR and PCR, most of these have focussed only on either of these tasks.

%Recently, researchers have tried to incorporate the idea of statutes during PCR in different ways, exploiting the inherent connection between the statutes and precedents cited from a given query, i.e., prior cases cited from a given query case are likely to involve same or similar statutes as those cited from the query case. Usually, these efforts have tried to incorporate the statutes cited from the precedent case, and used those as intermediate classification tasks. However, to the best of our knowledge, no work has previously tried to model the preidction of both the statutes and precedents cited from the same query simultaneously. In fact, there is no existing dataset that allows for the development of models to solve both tasks together. 

%The objective of this work to bridge the above gap, and to that end, we collected and curated a corpus consisting of Indian legal cases (from the Supreme Court and several High Courts of India) and Central Government laws, and created \ilpcsr{}, the first legal dataset that enables us to model LSR and PCR simultaneously from queries based on real-world cases. (see Table~\ref{tab:dataset_compare} for more details). Then, we performed several experiments on the dataset, including unsupervised lexical and embedding based approaches, supervised fine-tuning approaches, approaches to refine the input query text, etc. We used a variety of models like BM25, transformer-based models, GNNs as well as latest state-of-the-art LLMs as part of these approaches.

%Experiments and analyses on this first-of-its-kind dataset bring out some surprising observations -- lexical methods perform well in one task (PCR) while semantic methods perform well in the other (LSR). A probable reason for this surprising observation is the difference in the nature of the two tasks, such as, statutes are short and abstract, and the overall semantics of the query help to identify the relevant statutes. However, precedents are long and have similar language as the query (allows better lexical matches). Moreover, different portions of the query text (that have different local semantics) match accurately with different precedents. For this reason also, multi-tasking approaches also do not achieve much improvement.

% Finally, to overcome the above issues, we propose some simple ensemble-based approaches for combining scores coming from semantic and lexical methods respectively. These techniques use a linear interpolation of semantic and lexical scores, with a parameter $\alpha$ controlling the weightage between the two scores. We derive the $\alpha$ value by both grid-search and learning a linear layer over the semantic embeddings. Using the simple ensemble technique, we are able to obtain gains over both the individual methods. This observation holds true for both statute and precedent retrieval, across different combinations of lexical and semantic approaches. 
% We also conduct further experiments with the \coliee{} datasets, specifically \texttt{COLIEE 2024 Task 3} (Statute Retrieval on Japanese law) and \texttt{COLIEE 2024 Task 1} (Precedent Retrieval on Canadian Federal cases). We verify that the same trends as \ilpcsr{} hold for \coliee{}, and that the ensemble approach indeed provides an improvement here as well. 

% To sum up, our contributions are the following:

% \noindent (i)~We develop the first dataset that enables modeling of simultaneous retreival of both statutes and precedents cited from the same query.

% \noindent (ii)~We benchmark several approaches on both LSR and PCR tasks, and observe that while semantic features are more important for LSR, lexical features dominate PCR performance.

% \noindent (iii)~We develop simple ensemble based approaches to combine lexical and semantic approaches, and these ensembles perform better than either individual approach for both LSR and PCR across multiple datasets.
